\section{Curvature rank and section growth}
\label{sec:curvature-rank-section-growth}

The transcendence degree of a section-ratio field can be bounded by a smooth curvature form, even when that form changes rank and has no global sign. The key is to arrange that a maximum of a section norm occurs where the section map has its maximal rank. This gives a pointwise restriction on the growth counted by the preceding sections.

Keep the connected compact complex manifold $X$ and line bundle $L$, and choose a smooth Hermitian metric $h$ on $L$. In a local frame write $|e|_h^2=\mathrm e^{-\chi}$ and put $\beta=\ddc\chi$, as in \cref{chap:sections}. Let $r=\max_X\operatorname{rank}\beta$. No positivity of $\beta$, constant-rank condition or K\"ahler metric is assumed. Set $S_h=\{x:\beta_x\geq0\}$; when $S_h$ is nonempty, let $r_h^+=\max_{S_h}\operatorname{rank}\beta$.

\begin{proposition}\label{prop:section-map-curvature-rank}\index{Curvature rank!section maps}
For a nonzero subspace $W\subset H^0(X,L^p)$, $p\geq1$, let $B_W$ be its common zero set and let $\Phi_W\colon X\setminus B_W\to\mathbb P(W^*)$ be its section map. If $d$ is the maximal rank of its differential, there is $x_*\notin B_W$ such that
\[
 \operatorname{rank}\mathrm d\Phi_W|_{x_*}=d,
 \qquad \beta_{x_*}\geq0,
 \qquad \operatorname{rank}\beta_{x_*}\geq d.
\]
\end{proposition}
\begin{proof}
Choose a basis $s_0,\ldots,s_N$ of $W$. Principalizing the ideal of their local coefficients gives a smooth compact modification $\pi\colon Y\to X$, isomorphic over $X\setminus B_W$, and an effective divisor $D$ with
\[
 \pi^*s_j=s_Dv_j,\qquad
 v_j\in H^0(Y,\pi^*L^p(-D)),
\]
where the $v_j$ have no common zero. The resolved map $f=[v_0:\cdots:v_N]\colon Y\to\mathbb P^N$ has maximal differential rank $d$, since it agrees with $\Phi_W\circ\pi$ on a dense open set.

The critical set $C=\{y:\operatorname{rank}\mathrm df_y<d\}$ is analytic. Proper mapping and the image-dimension theorem imply that $Z=f(Y)$ and $T=f(C)$ are compact analytic sets, with $\dim Z=d$ and $\dim T\leq d-1$ when $T$ is nonempty. For the last inequality apply the rank bound to the smooth locus of every irreducible component of $C$. By Chow's theorem $T$ is algebraic. Choose a homogeneous polynomial $P$ of degree $a\geq1$ that vanishes on $T$ and does not vanish identically on $Z$. If $T$ is empty, take $P=1$ and $a=0$, which also covers $d=0$.

Then $\sigma=P(s_0,\ldots,s_N)$ is a nonzero section of $L^{pa}$. Indeed, the set where $P(v_0,\ldots,v_N)\ne0$ meets the dense locus where $\pi$ is an isomorphism. The smooth function
\[
 \Psi=\sum_j|\sigma s_j|_{h^{p(a+1)}}^2
\]
has a positive maximum $x_*$. At that point $\sigma\ne0$ and some $s_j\ne0$, so $x_*\notin B_W$ and its lift is outside $C$. The tuples $(\sigma s_j)$ and $(s_j)$ define the same projective map near $x_*$. The complex Hessian of $\log\Psi$ is seminegative at its maximum point $x_*$. With $\omega_{\mathrm{FS}}=\ddc\log(1+|w|^2)$ in an affine chart, this gives
\[
 0\geq\ddc\log\Psi|_{x_*}
   =\Phi_W^*\omega_{\mathrm{FS}}|_{x_*}
                    -p(a+1)\beta_{x_*}.
\]
The first form on the right is semipositive of rank $d$. Thus $\beta_{x_*}$ is semipositive and its kernel is contained in the kernel of that form, proving the assertion.
\end{proof}

If some positive power of $L$ has a nonzero section, the \emph{Kodaira--Iitaka dimension} $\kappa(L)$ is the transcendence degree of the section-ratio field of the full section ring; otherwise put $\kappa(L)=-\infty$.\index{Kodaira--Iitaka dimension}

\begin{corollary}\label{cor:curvature-rank-section-count}\index{Curvature rank!growth bound}
If $S_h$ is empty, then $H^0(X,L^p)=0$ for every $p\geq1$. Otherwise
\[
 \kappa(L)\leq r_h^+\leq r,
 \qquad h^0(X,L^p)=O(p^{r_h^+}).
\]
In this second case, every graded section algebra with a nonzero positive-degree piece has growth exponent at most $r_h^+$, and all its flag bodies, including those constructed on a point blowup, have dimension at most $r_h^+$. For every fixed weight $u\in\QPSH(X)$, the integrable spaces satisfy $\dim\mathcal H_p(u)=O(p^{r_h^+})$.
\end{corollary}
\begin{proof}
The proposition applied to a single section excludes every positive power if $S_h$ is empty. If all positive-degree section spaces vanish, the remaining section bounds are immediate and $\kappa(L)=-\infty$. Otherwise, suppose a graded algebra has a nonzero positive-degree piece. Write finitely many algebraically independent elements of its ratio field with a common denominator in one positive degree, as in the proof of \cref{prop:flag-dimension-section-field}. They are ratios of coordinates of the section map of the span $W$ of the numerators and the denominator. The image of the resolved section map is a projective variety, by proper mapping and Chow's theorem, and the given elements are pullbacks of rational functions on it; its dimension is therefore at least their number. This dimension is the maximal rank of $\mathrm d\Phi_W$, which the proposition bounds by $r_h^+$. Hence the transcendence degree of the ratio field is at most $r_h^+$.

Apply \cref{thm:all-compact-section-count} to this algebra and to the full section ring. The degree-period formula gives the growth bound in every degree, including the zero spaces outside the period. By \cref{thm:intrinsic-flag-count}, the dimension of every flag body equals the growth exponent. Finally, $\mathcal H_p(u)\subseteq H^0(X,L^p)$, so the same upper bound holds even though the integrable spaces need not form a graded algebra.
\end{proof}

\begin{lemma}\label{lem:curvature-rank-fixed-line-twist}
For every fixed holomorphic line bundle $A$ on $X$,
\[
 h^0(X,L^p\otimes A)=O(p^r).
\]
\end{lemma}
\begin{proof}
Induct on the dimension, simultaneously for all connected compact complex manifolds and all integer upper bounds on curvature rank. Dimension zero is immediate. If $H^0(X,L^p\otimes A)=0$ for every $p\geq1$, there is nothing to prove. Otherwise fix a nonzero section of $L^a\otimes A$ for some $a\geq1$, and resolve its zero divisor on a smooth compact modification $\mu\colon Y\to X$. With $M=\mu^*L$, the pulled-back twist is $M^{-a}(D)$ for a fixed effective divisor $D$ whose components are smooth. Pullback is injective on sections, and the pulled-back curvature has rank at most $r$.

Write $D=H_1+\cdots+H_s$, allowing repeated components, and put $D_j=H_1+\cdots+H_j$. For $m=p-a>0$, the divisor sequences give
\[
 h^0(Y,M^m(D))
 \leq h^0(Y,M^m)+\sum_{j=1}^s
       h^0\bigl(H_j,(M|_{H_j})^m\otimes\mathcal O_Y(D_j)|_{H_j}\bigr).
\]
The first term is $O(m^r)$ by the preceding corollary; every other term is $O(m^r)$ by induction, since curvature rank cannot increase on restriction. The number of terms and their twists are fixed. Absorb the finitely many $p\leq a$ in the constant.
\end{proof}

\begin{corollary}\label{cor:curvature-rank-endpoint-cohomology}\index{Dolbeault cohomology!curvature rank at endpoints}
For every fixed holomorphic vector bundle $E$ on $X$,
\[
 h^0(X,L^p\otimes E)=O(p^r),
 \qquad h^n(X,L^p\otimes E)=O(p^r).
\]
\end{corollary}
\begin{proof}
For $E$ of positive rank, let $\pi\colon Y=\mathbb P(E^*)\to X$ be the bundle of lines in $E^*$, and let $\mathcal O_Y(1)$ be the dual of its tautological line bundle. Fiberwise linear forms give $\pi_*\mathcal O_Y(1)=E$, hence
\[
 H^0(X,L^p\otimes E)
 =H^0(Y,(\pi^*L)^p\otimes\mathcal O_Y(1)).
\]
The pulled-back curvature has rank at most $r$, so the fixed-line-twist lemma, proved in all dimensions, gives the first bound. Serre duality identifies the dual of the top cohomology with
\[
 H^0(X,(L^*)^p\otimes E^*\otimes K_X).
\]
The dual curvature is $-\beta$ and has the same rank; apply the bound for $h^0$ to $L^*$, with $E^*\otimes K_X$ in place of $E$. A bundle of rank zero has zero cohomology and satisfies both bounds.
\end{proof}

\begin{proposition}\label{prop:curvature-rank-euler-polynomial}\index{Euler characteristic!curvature rank}
For every fixed holomorphic vector bundle $E$, the Euler characteristic $\chi(X,L^p\otimes E)$ is a polynomial in $p$ of degree at most $r$.
\end{proposition}
\begin{proof}
The Hirzebruch--Riemann--Roch theorem, which holds on every compact complex manifold, gives
\[
 \chi(X,L^p\otimes E)
 =\int_X\mathrm e^{p\beta}\operatorname{ch}(E)
                              \operatorname{td}(T^{1,0}X)
 =\sum_{j=0}^{r}\frac{p^j}{j!}
       \int_X\beta^j\wedge\operatorname{ch}(E)
                              \wedge\operatorname{td}(T^{1,0}X).
\]
Here the integral selects real degree $2n$, and the Chern character and Todd forms are computed with fixed Hermitian metrics. The curvature normalization represents $c_1(L)$, and $\beta^{r+1}=0$ pointwise by its rank bound. Every coefficient is independent of $p$.
\end{proof}

\begin{corollary}\label{cor:surface-curvature-rank-cohomology}
If $\dim X=2$, then $h^q(X,L^p\otimes E)=O(p^r)$ for every fixed holomorphic vector bundle $E$ and $0\leq q\leq2$.
\end{corollary}
\begin{proof}
The endpoints have already been bounded. Since $\chi=h^0-h^1+h^2$, the middle dimension is at most $h^0+h^2+|\chi|=O(p^r)$ by the Euler polynomial.
\end{proof}

\begin{example}\label{ex:curvature-rank-sharp-section-growth}
Let $0\leq r\leq n$, let $Y$ be any connected compact complex manifold of dimension $n-r$, and take $X=\mathbb P^r\times Y$ and $L=\operatorname{pr}_1^*\mathcal O(1)$ with the pulled-back Fubini--Study metric. Its curvature has rank $r$, and
\[
 h^0(X,L^p)=\binom{p+r}{r}.
\]
Thus the section-growth exponent cannot be reduced, including when $Y$ is non-K\"ahler.
\end{example}
\begin{proof}
A section, expressed in a line-bundle frame on the projective factor, is a holomorphic function of $Y$ and hence is constant there. Restriction to any slice $\mathbb P^r\times\{y_0\}$ is therefore inverse to pullback on sections. Counting degree-$p$ homogeneous monomials in $r+1$ variables gives the formula; for $r=0$ it equals one.
\end{proof}

These bounds concern the closed curvature of a line bundle. They do not estimate the mass of a potential in an unrelated nonclosed Hermitian background, as \cref{prop:sections-infinite-mass-obstruction} shows. In dimension at least three, the endpoint bounds and the Euler polynomial alone do not control the individual intermediate cohomology dimensions. The Hopf body computed next realizes the rank bound $n-1$ for the pulled-back positive line bundles of \cref{prop:hopf-flag-simplex}.
